A right circular cone of radius $5$ and height $12$ is laid on the ground on its side, then pushed so that it turns around in a circle, rolling without slipping about its vertex. After the cone's base has drawn out a complete circle on the ground, the cone has completed some number of spins about its axis. This number can be written as a fraction $\dfrac{m}{n}$ where $m$ and $n$ are relatively prime positive integers. What is $m + n$?
